\subsection{The two markings of the descent map}

\paragraph{A4.M1 (Proved): the literal map in the current marking is a different map.}
Let $\widehat B=\langle s,r\mid s^2=r^3=1\rangle^{\wedge}$,
$\tau=sr$, and let $F$ be its level-two subgroup.  Distinguish the
two free generating pairs
\[
 x_P=\tau^{-2},\quad y_P=sx_Ps=(rs)^{-2},\qquad
 x_L=(rs)^2=y_P^{-1},\quad y_L=\tau^2=x_P^{-1}.
\]
The subscripts $P,L$ mean the current prompt and LM2, respectively.
The original manuscript/LM2 quotient, expressed on the physical subgroup
$F$, is
\[
 \varphi_L(x_L)=s,\quad\varphi_L(y_L)=\tau,
 \qquad\text{hence}\qquad
 \varphi_L(x_P)=\tau^{-1},\quad\varphi_L(y_P)=s.
\]
The literal quotient specified in the current prompt is instead
\[
 \varphi_P(x_P)=s,\qquad\varphi_P(y_P)=\tau.
\]
Let $\kappa\in\operatorname{Aut}(F)$ send $x_P$ to $y_P^{-1}$ and
$y_P$ to $x_P^{-1}$. Then $\kappa^2=1$ and
\[
                 \varphi_L=\varphi_P\circ\kappa.                 \tag{M1}
\]
In particular, retaining the written values $x\mapsto s,y\mapsto\tau$
while changing from the LM2 basis to the current basis does not retain
the same physical homomorphism.

\emph{Proof.} All assertions follow by substituting the four displayed
words in $s,r$. The automorphism $\kappa$ is $\sigma\circ c$, where
$\sigma$ exchanges $x_P,y_P$ and $c$ inverts both generators.
The two factors commute.  Equation (M1) can be checked on the two
free generators.\hfill$\square$

\paragraph{A4.M2 (Proved): the two descent endomorphisms differ by conjugation by $c$.}
Let $G=\widehat{\mathrm{GT}}_0$ be the outer centralizer of the deck
$S_3$ acting on $F$. Denote by $D_P,D_L$ the endomorphisms obtained by
outer descent along $\varphi_P,\varphi_L$, followed by the normalized
restriction to $F$. Then
\[
 D_L(g)=D_P(cgc)=cD_P(g)c,\qquad
 D_P(c)=c.                                                       \tag{M2}
\]
Thus the two compatibility subgroups are
\[
 C_L=\operatorname{Fix}(D_L)
     =\{g:D_P(g)=cgc\},\qquad
 C_P=\operatorname{Fix}(D_P).
\]
No equality of these two subgroups is asserted.

\emph{Proof.} For any quotient $\varphi$ and an automorphism $a$
preserving its kernel, write $\operatorname{Desc}_{\varphi}(a)$ for
the induced outer automorphism of its target. Equation (M1) gives
\[
 \operatorname{Desc}_{\varphi_L}(g)
 =\operatorname{Desc}_{\varphi_P}(\kappa g\kappa^{-1}).
\]
In the outer automorphism group, $g$ commutes with $\sigma$, while
$\kappa=\sigma c=c\sigma$. Therefore
$\kappa g\kappa^{-1}=cgc$. Applying normalized restriction proves
the first equality in (M2).

For $c(x_P)=x_P^{-1},c(y_P)=y_P^{-1}$, descent along $\varphi_P$
is the genuine automorphism
\[
 \iota(s)=s,\qquad\iota(\tau)=\tau^{-1},\qquad
 \iota(r)=sr^{-1}s.
\]
It induces the identity on $B/F=S_3$, since conjugation by the
transposition $s$ inverts the three-cycle. Its restriction sends
$x_P$ and $y_P$ to their inverses. Hence $D_P(c)=c$.
The homomorphism law for $D_P$ now proves the second equality in (M2).
All statements are about outer maps; inner choices disappear before
normalized restriction.\hfill$\square$

\paragraph{A4.M3 (Proved): the original arithmetic compatibility and a counterexample to the literal-current fixed-point assertion.}
In the original LM2 marking, arithmetic functoriality gives
$D_L(\sigma)=\sigma$ for every $\sigma\in G_{\mathbb Q}$.
Consequently the literal-current endomorphism satisfies
\[
                   D_P(\sigma)=c\sigma c
                   \quad(\sigma\in G_{\mathbb Q}).              \tag{M3}
\]
It does not fix every Galois element. In fact,
\[
                  C_P\cap G_{\mathbb Q}=\{1,c\}.
\]

\emph{Proof of the arithmetic comparison.} Let $\mathcal X$ be the
orbifold with coarse curve $\mathbb P^1_{\mathbb Q}\setminus\{0\}$,
with inertia orders two at $1$ and three at $\infty$. It is the
rigidified elliptic modular stack in coordinate $T=1728/j$.
Its level-two cover and its puncture-opening map are, respectively,
\[
 U\longrightarrow\mathcal X,\quad
 T=\frac{27t^2(1-t)^2}{4(1-t+t^2)^3},
 \qquad
 U=\mathbb P^1\setminus\{0,1,\infty\}\hookrightarrow\mathcal X.
\]
Here is a simultaneous choice of loops, including their orientation.
Choose a small rational $t_0>0$, for instance $t_0=1/100$, and base
the orbifold at $T_0=q(t_0)$, where $q$ is the displayed rational
function. Let $\tau$ be the positive, counterclockwise meridian
about $T=0$, and let $s$ be the positive meridian about $T=1$,
using stems along the real interval, slightly displaced when needed.
With the peripheral product ordered as $\ell_1\ell_0\ell_\infty=1$,
the positive infinity meridian is $(s\tau)^{-1}$. Thus setting
$r=s\tau$ gives $s^2=r^3=1$ and $\tau=sr$.

The cover has the deck involution $t\mapsto1-t$ and local degree
two at $t=0,1,\infty$. Its only other ramifications have degree
two over $T=1$ and degree three over $T=\infty$; explicitly
\[
 1-q(t)=\frac{(t+1)^2(2t-1)^2(t-2)^2}
                  {4(t^2-t+1)^3}.
\]
Together with the six anharmonic deck transformations this shows
that it is the degree-six $S_3$ orbifold cover.
Let $y_L$ be the positive meridian about $t=0$ based at $t_0$.
Since $q(t)=(27/4)t^2+O(t^3)$, its image is $\tau^2$, with the
positive exponent. The lift $\gamma$ of the positive loop $s$
starting at $t_0$ ends at $1-t_0$: it passes around the ramification
point $1/2$ and implements the deck involution $t\mapsto1-t$.
Define $x_L$ by transporting the positive meridian about $t=1$
back along $\gamma$. The latter meridian is the image of $y_L$
under the holomorphic involution, so it has the same positive
orientation. Hence
\[
 q_*(y_L)=\tau^2=(sr)^2,\qquad
 q_*(x_L)=s\tau^2s^{-1}=(rs)^2.
\]
The stems can be chosen with disjoint interiors, so $x_L,y_L$ are
a free peripheral basis of $\pi_1(U,t_0)$. These formulas therefore
identify $q_*$ with the physical level-two inclusion in precisely
the LM2 marking, rather than merely an unspecified conjugate or
orientation reversal of that inclusion.

For the open map use the real interval from $T_0$ to $t_0$ to
transport its basepoint. In the target orbifold the path $\gamma$
is homotopic to the interval from $t_0$ to $1-t_0$: its detour is
around $1/2$, which is an ordinary interior point of the target.
Thus the same two domain loops have images $\tau$ and $s$ under
the open map, respectively. This proves simultaneously
\[
 (x_L,y_L)\longmapsto((rs)^2,(sr)^2)\quad\text{under }q_*;
 \qquad (x_L,y_L)\longmapsto(s,\tau)\quad\text{under }j_*.
\]
Both algebraic maps are defined over $\mathbb Q$. The interval
used to transport basepoints need not be Galois invariant: changing
its Galois transform changes the induced comparison by a geometric
inner automorphism, exactly the ambiguity already removed in the
outer descent and restriction. Functoriality of the arithmetic
fundamental-group extensions therefore makes these two outer
actions equivariant. Consequently $D_L$ fixes the Galois image.
This reproves the typed content of LM2 E.1 with the same physical
marking on both maps; it does not prove the same assertion with
the unconverted formula for $\varphi_P$.
There is one final identification of arithmetic images to check.
Changing the positive LM2 basis to the prompt basis conjugates the
outer Galois image by $\kappa$. Since $\kappa=\sigma c$, the
interchange $\sigma$ commutes with this image in the outer group,
and $c$ belongs to the image, one has
$\kappa G_{\mathbb Q}\kappa^{-1}=G_{\mathbb Q}$ as subgroups.
Thus the pointwise fixed subgroup statement for $D_L$ applies to
every element of the current Galois subgroup, without identifying
the old and new labels of individual arithmetic automorphisms.
Combining it with (M2) proves (M3).

For completeness the needed Galois centralizer fact has a direct
field-theoretic proof. An automorphism commuting with $c$ preserves
its fixed field $\mathbb R_{\rm alg}=\overline{\mathbb Q}\cap\mathbb R$.
The positive elements of this field are exactly its nonzero squares,
so every field automorphism preserves its order. It fixes the dense
subfield $\mathbb Q$, hence fixes every real algebraic number: two
different real numbers are separated by a rational number. Finally
$\overline{\mathbb Q}=\mathbb R_{\rm alg}(i)$, so an automorphism
fixing $\mathbb R_{\rm alg}$ is either $1$ or $c$. Therefore
$C_{G_{\mathbb Q}}(c)=\{1,c\}$, proving the displayed intersection
without a GT centralizer theorem.\hfill$\square$

Here is an explicit compact family of witnesses that does not rely
only on the centralizer theorem. Choose compatible positive real
roots $2^{1/n}$ and standard roots of unity $\zeta_n$. Define the
Kummer cocycle by
\[
 \sigma(2^{1/n})/2^{1/n}=\zeta_n^{\rho_2(\sigma)}.
\]
The number $2^{1/4}$ does not belong to $\mathbb Q^{\mathrm{ab}}$:
$X^4-2$ is Eisenstein, and its real root generates a nonnormal real
quartic field, whereas every finite subextension of
$\mathbb Q^{\mathrm{ab}}/\mathbb Q$ is Galois. Its square does belong
to $\mathbb Q^{\mathrm{ab}}$. Thus
\[
 \mathcal C=
 \{\sigma\in G_{\mathbb Q^{\mathrm{ab}}}:
                    \sigma(2^{1/4})=-2^{1/4}\}
\]
is a nonempty compact coset, and each of its elements has
$\rho_2(\sigma)\equiv2\pmod4$.
Complex conjugation fixes every chosen positive root and inverts
every $\zeta_n$, so
\[
                 \rho_2(c\sigma c)=-\rho_2(\sigma).
\]
In $\widehat{\mathbb Z}$, a nonzero element cannot equal its negative:
multiplication by two is injective. More concretely, a residue equal
to two modulo four is two or six modulo eight, and changes under
negation. Every $\sigma\in\mathcal C$ therefore satisfies
$D_P(\sigma)=c\sigma c\ne\sigma$. The claimed literal-current
Galois fixed-point assertion is refuted by this entire compact coset.
The same witnesses refute both
$\widehat{\mathrm{GT}}\subseteq C_P$ and $C_P=\ker\Delta$:
they belong to $G_{\mathbb Q}\subseteq\widehat{\mathrm{GT}}$ and
have $\Delta=0$ by R4, but do not belong to $C_P$. These
counterexamples do not refute the corresponding proposed inclusion
for the original group $C_L$.

\paragraph{A4.M4 (Proved): the exact multiplier-one formula in the current marking.}
Write $x=x_P$, $y=y_P$, $z=(xy)^{-1}$ and
$\sigma=\operatorname{Ad}(s)|_F$, $\rho=\operatorname{Ad}(r)|_F$.
For a normalized pair $(1,f)$, set
\[
 C=\varphi_P(f)=f(s,\tau),\qquad
 H=\varphi_P(f(z,x))=f(r^{-1},s).
\]
Then $C\in F$. Write $C^{\mathrm{ab}}=(p,q)$ in the basis $[x],[y]$,
put $a=q-p$, and write $H\bmod F=r^l$ with $l\in\mathbb Z/3$.
The normalized pair and the two pure frames of $D_P(1,f)$ are
\[
 \begin{split}
 f_{D_P}&=y^a\sigma(C)C^{-1}x^{-a},\\
 \beta_{D_P}&=x^a C,\\
 A_{D_P}&=x^a CHr^{-l}.
 \end{split}                                                    \tag{M4}
\]
In particular, the current literal fixed-point condition is exactly
\[
             f=y^{q-p}\sigma(\varphi_P(f))\varphi_P(f)^{-1}
                    x^{p-q}.                                   \tag{M5}
\]
The corresponding original compatibility condition is
$D_P(1,f)=c(1,f)c$, equivalently
$f_{D_P}=f(x^{-1},y^{-1})$.

\emph{Proof.} The standard multiplier-one representative fixes $x$
and sends $y$ to $f^{-1}yf$. Its extension to $B$ is
$\Phi(s)=sf$, $\Phi(r)=f^{-1}r$, $\Phi(\tau)=\tau$.
The three-cycle relation gives
\[
 \Phi(z)=f^{-1}\rho(f)^{-1}z\rho(f)f
        =f(z,x)z f(z,x)^{-1}.
\]
Thus its descendant $\psi$ satisfies
\[
 \psi(s)=s,\qquad\psi(\tau)=C^{-1}\tau C,
 \qquad\psi(r)=HrH^{-1}.
\]
Because $f\in F'$, both $C,H$ map into the derived subgroup $A_3$
of $S_3$. The last formula therefore fixes $r$ modulo $F$, and the
first fixes $s$. Hence $\psi$ induces the identity on $S_3$.
The middle formula now shows that $C\bmod F$ centralizes the
transposition $\tau$, and $A_3\cap C_{S_3}(\tau)=1$, proving $C\in F$.

The map $\operatorname{Ad}(C)\psi$ fixes $\tau$ and has pure frames
$C$ and $CHr^{-l}$. Its pair coordinate is
$\sigma(C)C^{-1}$, with abelianization $(q-p,p-q)=(a,-a)$.
Composing by $\operatorname{Ad}(x^a)$ fixes the cusp and sends this
coordinate to $y^a\sigma(C)C^{-1}x^{-a}\in F'$. This proves (M4),
including the normalization exponent. Since both normalized pairs
fix $x$ and have coordinate in $F'$, they represent the same outer
class only when their coordinates are equal; their possible inner
conjugator centralizes $x$ and its abelianization forces its exponent
to vanish. This proves (M5).\hfill$\square$

\paragraph{A4.M5 (Proved): an independent check from the NT formulas.}
The primary report of Tsunogai and Nakamura (2002, Theorems 1 and 2)
uses braid generators $t_1,t_2$ and $\eta=t_1t_2t_1$.
The exact quotient to the current $B$ marking is
\[
 t_1=\tau^{-1}=r^{-1}s,\qquad t_2=s\tau^{-1}s=sr^{-1},
 \qquad t_1^2=x,\ t_2^2=y,\ \eta=s,\ t_1t_2=r.
\]
Let $\gamma\in G_{\mathbb Q^{\mathrm{ab}}}$, and write its harmonic
primitive as $g=\beta^{-1}$. The source gives
\[
 g=f(x,s)t_1^{4\rho_2},\qquad
 f(t_1,s)=s^{\rho_3-2\rho_2}f(x,s)t_1^{6\rho_2-3\rho_3}.
\]
Apply these formulas to $c\gamma c$. Its pair coordinate is
$f^c(X,Y)=f(X^{-1},Y^{-1})$, its Kummer exponents are
$-\rho_2,-\rho_3$, and its pure harmonic frame is $\iota(\beta)$,
where $\iota$ is the reflection used in A4.M2. As both Kummer
exponents of $\gamma$ are even, the order-two factor disappears.
Taking the inverse yields
\[
 C=f(s,\tau)
    =t_1^{2\rho_2-3\rho_3}\,\iota(\beta).
\]
Since $\beta^{\mathrm{ab}}=(-\rho_2,-\rho_2)$, the second
coordinate of $C^{\mathrm{ab}}$ is $q=\rho_2$.
Formula (M4) gives $\rho_2(D_P\gamma)=-q=-\rho_2(\gamma)$,
in agreement with (M3), and its normalization makes
$\beta_{D_P\gamma}=\iota(\beta_\gamma)$ exactly.
This independently detects the conjugation twist using the opened
primary equations, rather than a presumed identification with them.

\paragraph{Consequences for the unresolved comparison.}
The manuscript equation (30) remains correct in its original marking:
with $\xi=y_L^{(\lambda-1)/2}f_L$ and an allowed gauge,
\[
 \xi\rho(\xi)\rho^2(\xi)=1,\qquad
 \xi=[\varphi_L(\xi),s]x_L^{(\lambda-1)/2}.
\]
Its literal unconverted use with $x_P,y_P,\varphi_P$ changes the
compatibility subgroup. This correction does not by itself prove
$\widehat{\mathrm{GT}}\subseteq C_L$ or $C_L\subseteq\ker\Delta$.
For the finite-character route, however, the distinction is harmless
at the level of the composed character: because $\Delta$ is invariant
under conjugation by $c$,
\[
                         \Delta\circ D_L=\Delta\circ D_P.
\]
A proof that this common character vanishes would put both fixed
subgroups in $\ker\Delta$. No such vanishing is presumed here.

\paragraph{A4.M6 (Proved): an exact NT primitive comparison and two conditional inclusions.}
Let $\Phi=(1,f)\in\mathcal H$, and put
\[
 a=\tau^{-1},\quad b=s\tau^{-1}s,\quad a^2=x,\quad b^2=y,
 \quad aba=s,\quad ab=r.
\]
Write $g=\beta^{-1}$, $h=A^{-1}$, so that
\[
 f=\sigma(g)^{-1}g=\rho(h)^{-1}h,
 \qquad g^{\rm ab}=(\rho_2,\rho_2),\quad h^{\rm ab}=0.
\]
The last equality follows from the order-three primitive equation:
the matrix $1-\rho$ on $F^{\rm ab}$ has determinant three and is
injective over $\widehat{\mathbb Z}$. The coordinate $\rho_2$ is even
on the whole multiplier-one group, as established in the supplied
R4 derivation.

Set $P=f(a,s)$. Then $P\in F$. Write $P^{\rm ab}=(p,q)$ and
$b_0=p-q$. Descent along the \emph{original physical} quotient
$\varphi_L(x)=a,\varphi_L(y)=s$ gives
\[
 \beta_{D_L}=x^{b_0}P^{-1}.                                     \tag{M6}
\]
Consequently
\[
 \Phi\in C_L\quad\Longleftrightarrow\quad
       \beta^{-1}=f(a,s)x^{-b_0}.                              \tag{M7}
\]
In particular, the current fixed-point dictionary identifies the
original compatibility with the \emph{GF0 primitive form}. It does
not directly identify it with the harmonic GF1 form.

More precisely, fix any even $\kappa\in\widehat{\mathbb Z}$ and consider
the four equations in $\widehat B$
\[
\begin{aligned}
 (\mathrm{GF}_0)_\kappa:\quad
       g&=f(a,s)a^{-2\rho_2+3\kappa},\\
 (\mathrm{GF}_1):\quad
       g&=f(x,s)a^{4\rho_2},\\
 (\mathrm{HF}_0)_\kappa:\quad
       h&=r^{\rho_2-3\kappa/2}f(a,r)a^{3\kappa-2\rho_2},\\
 (\mathrm{HF}_1)_\kappa:\quad
       h&=f(x,r)a^{3\kappa}.
\end{aligned}                                                   \tag{M8}
\]
These are the central-quotient forms of NT Theorem 2 at multiplier
one, with $\kappa$ in place of its Kummer parameter $\rho_3$.
Then
\[
 \begin{split}
 \Phi\in C_L
 &\Longleftrightarrow
   (\mathrm{GF}_0)_\kappa\text{ for some even }\kappa\\
 &\Longleftrightarrow
   (\mathrm{HF}_0)_\kappa\text{ for some even }\kappa.
 \end{split}                                                    \tag{M9}
\]
The parameter is unique when either equation holds. Moreover the
following two sufficient conditions for $C_L$ are proved:
\begin{align}
 &(\mathrm{GF}_1)\quad\text{and}\quad
 f(a,s)=f(x,s)a^{6\rho_2-3\kappa};                              \tag{M10}\\
 &(\mathrm{HF}_1)_\kappa\quad\text{and}\quad
 f(a,r)=r^{-\rho_2}f(x,r)a^{2\rho_2}.                           \tag{M11}
\end{align}
The second equations in (M10) and (M11) are respectively the
central-quotient forms of NT (V) and $(\mathrm{IV}'_{\rm bis})$.
NT's harmonic equation $(\mathrm I')$ is equivalent to
$(\mathrm{GF}_1)$, and its equianharmonic equation $(\mathrm{II}')$
is equivalent to $(\mathrm{HF}_1)$ with the stated parameters.
Thus (harmonic $+$ V) and (equianharmonic $+$ IV-bis) each give a
precisely typed inclusion in the original compatibility group.
This is an inclusion between explicitly defined equation loci;
it is not a claim that all of GT satisfies those additional equations.

\emph{Proof.} The descendant along $\varphi_L$ fixes $a$ and sends
$s$ to $P^{-1}sP$. Its image of $r$ is $QrQ^{-1}$, where
$Q=f(a,r)^{-1}$. Since $P,Q$ are commutators in $B$, their reductions
lie in $A_3$; the descendant fixes $r$ and the transposition $a$
modulo $F$, hence fixes all of $S_3$. This proves $P\in F$.
The initial $s$-frame is $P^{-1}$ and has abelianization $(-p,-q)$.
The unique normalization making the pair coordinate a commutator
is conjugation by $x^{p-q}$. This proves (M6). The normalized map
fixes $\tau$, so its $s$-frame determines the whole automorphism;
this proves (M7).

Suppose first that $\Phi\in C_L$. Equation (M7) gives $q=\rho_2$.
As $P\in B'$, its image in $B^{\rm ab}$ is zero. The images of
$x,y$ in the order-three factor of $B^{\rm ab}$ coincide, so
$p+q\equiv0\pmod3$. Hence
\[
      \kappa=\frac{2(\rho_2-b_0)}3
\]
is well defined and even, and (M7) is exactly
$(\mathrm{GF}_0)_\kappa$. Conversely this primitive equation makes
$P^{\rm ab}=(2\rho_2-3\kappa/2,\rho_2)$, so
$b_0=\rho_2-3\kappa/2$ and (M7) follows. The cusp subgroup is
$\widehat{\mathbb Z}$, on which multiplication by three is injective,
so $\kappa$ is unique.

For the order-three frame write $Q\bmod F=r^j$.
The normalized descendant has
\[
 A_{D_L}=x^{b_0}Qr^{-j}.
\]
Its abelianization in $F$ is zero. Passing further to $B^{\rm ab}$
therefore gives $j\equiv b_0\pmod3$. Inverting this formula yields
$h_{D_L}=r^{b_0}f(a,r)x^{-b_0}$. It is
$(\mathrm{HF}_0)_\kappa$ with the same $\kappa$ when $\Phi\in C_L$.
Conversely, if $(\mathrm{HF}_0)_\kappa$ holds, its inverse is
$A=x^dQr^{-d}$ with $d=\rho_2-3\kappa/2$. Purity gives
$d\equiv j\pmod3$, and comparison with $A_{D_L}$ makes their
quotient $x^{d-b_0}$. Both frames have zero abelianization, so
$d=b_0$. The two normalized maps fix $\tau$ and have the same
$r$-frame, hence are equal. This proves (M9).

Combining the two equations in (M10) gives
$g=f(a,s)a^{-2\rho_2+3\kappa}$ directly. Combining those in (M11)
gives
\[
 r^{\rho_2-3\kappa/2}f(a,r)a^{3\kappa-2\rho_2}
 =r^{-3\kappa/2}f(x,r)a^{3\kappa}
 =f(x,r)a^{3\kappa}=h,
\]
where the middle equality uses $r^3=1$ and evenness of $\kappa$.
Finally, inserting the GF1 or HF1 expression in the corresponding
primitive equation produces NT $(\mathrm I')$ or $(\mathrm{II}')$;
the reverse implication follows from uniqueness of the pure
primitive. All substitutions use $a^2=x,b^2=y,aba=s,ab=r$, so
there is no unrecorded change of braid generators.\hfill$\square$

\paragraph{What this says about (III$'$).}
Nakamura--Schneps (2000, Theorem 1.1) puts its relation (III$'$) in
the five-point mapping class group, with
$x_{45}=(t_1t_2)^3$ and $x_{51}=(t_2t_3)^3$; Proposition 5.4
gives the central correction when lifting it to $B_4$.
Tsunogai--Nakamura (2002, p.200) separately displays
$(\mathrm I'),(\mathrm{II}'),(\mathrm{III}'),(\mathrm{IV}')$ and (V).
The opened texts do not state that (III$'$) alone supplies (V), or
that all GT elements satisfy GF0. Equations (M8)--(M11) specify
the remaining comparison exactly: to pass from the harmonic
primitive to the original $C_L$ condition one needs the displayed
V comparison (or an independently proved substitute for it).
